% ----------------------------------------------------------------------
%              Latex PhD template for the University of Deusto
% ----------------------------------------------------------------------


%: Style file for Latex
% Most style definitions are in the external file PhDthesisPSnPDF.
% In this template package, it can be found in ./Latex/Classes/
\documentclass[twoside,12pt]{Latex/Classes/PhDthesisPSnPDF}

% My own packages
%\usepackage{pdflscape}
%\usepackage{longtable}
\usepackage{algpseudocode}
\usepackage{minted}
\usepackage[printonlyused]{acronym}
\usepackage{enumerate}
\usepackage{footnote} % footnotes inside tables using savenotes
\usepackage[inline]{enumitem} % inline lists
% For inline cites (section 6.2
\usepackage{bibentry} % http://gking.harvard.edu/files/bibentry2.pdf
%\nobibliography* throws an error
\nobibliography{9_backmatter/ref_self} % http://stefaanlippens.net/bibentry


%: Macro file for Latex
% Macros help you summarise frequently repeated Latex commands.
% Here, they are placed in an external file /Latex/Macros/MacroFile1.tex
% An macro that you may use frequently is the figuremacro (see introduction.tex)
%---------------------------------------------------------------
% Macros
% version 3 by Igor Ruiz-Agundez 2011
% version 2 by Jakob Suckale 2007
% version 1 by Harish Bhanderi 2002
%---------------------------------------------------------------

% This file contains macros that can be called up from connected TeX files
% It helps to summarise repeated code, e.g. figure insertion (see below).



%---------------------------------------------------------------
% Figures
%---------------------------------------------------------------


% Makes the \InsertFig macro compatible both with one or two columns
\makeatletter
\newlength \figwidth
\if@twocolumn
  \setlength \figwidth {\columnwidth}
\else
  \setlength \figwidth {\textwidth}
\fi
\makeatother

% \InsertFig allows inserting figures
% Parameters
% 1 --> Filename
% 2 --> Label for referencing
% 3 --> Title describing the figure (caption)
% 4 --> Description of the figure
% 5 --> Figure width, range [0,1]. If parameter is left blank the figure size is not change
% 6 --> Any other option for \includegraphics
% Usage:
% \InsertFig{}{}{}{}{}{}
%

\newcommand{\InsertTab}[5]{
  \ifthenelse{\isempty{#5}}
  { % if empty
    \begin{savenotes}
    \begin{table}
      \caption{#2}{\footnotesize{#3} \medskip}
      
      \centering
      #4
      \label{#1}
    \end{table}
    \end{savenotes}
  }
  { % if not empty
    \begin{savenotes}
    \begin{table}[#5]
      \caption{#2}{\footnotesize{#3} \medskip}
      
      \centering
      #4
      \label{#1}
    \end{table}
    \end{savenotes}
  }
}

\newcommand{\InsertFig}[6]{%
	\ifthenelse{\isempty{#5}}%
	{% if #1 is empty
		\begin{figure}[htbp!]
		\centering
		\includegraphics[#6]{#1}%
		\caption{#3}{\footnotesize{#4}}
		\label{#2}
		\end{figure}    
	}
	{% if #1 is not empty
		\begin{figure}[htbp!]
		\centering
		\includegraphics[width=#5\figwidth,#6]{#1}%
		\caption{#3}{\footnotesize{#4}}
		\label{#2}
		\end{figure}
	}
}

%% Simple version of \InsertFig
%\newcommand{\InsertFig}[5]{
%  \begin{figure}[htbp]
%   	\centering
%    \includegraphics[width=#4\textwidth,#5]{#1}%
%    \caption{#3}
%    \label{#2}
%  \end{figure}
%}



% insert a centered figure with caption
% parameters 1:filename, 2:label, 3:title, 
\newcommand{\figuremacro}[3]{
	\begin{figure}[htbp]
		\centering
		\includegraphics[width=1\textwidth]{#1}
		\caption[#3]{\textbf{#3}}
		\label{#2}
	\end{figure}
}


% insert a centered figure with caption and description
% parameters 1:filename, 2:label, 3:title, 4:description
\newcommand{\figuremacroD}[4]{
	\begin{figure}[htbp]
		\centering
		\includegraphics[width=1\textwidth]{#1}
		\caption[#3]{\textbf{#3} - #4}
		\label{#2}
	\end{figure}
}

% insert a centered figure with caption and description AND WIDTH
% parameters 1:filename, 2:label, 3:title, 4:description, 5: textwidth
% textwidth 1 means as text, 0.5 means half the width of the text
\newcommand{\figuremacroDW}[5]{
	\begin{figure}[htbp]
		\centering
		\includegraphics[width=#5\textwidth]{#1}
		\caption[#3]{\textbf{#3} - #4}
		\label{#2}
	\end{figure}
}

% inserts a figure with wrapped around text; only suitable for NARROW figs
% o is for outside on a double paged document; others: l, r, i(inside)
% text and figure will each be half of the document width
% note: long captions often crash with adjacent content; take care
% in general: above 2 macro produce more reliable layout
\newcommand{\figuremacroN}[3]{
	\begin{wrapfigure}{o}{0.5\textwidth}
		\centering
		\includegraphics[width=0.48\textwidth]{#1}
		\caption[#2]{{\small\textbf{#2} - #3}}
		\label{#1}
	\end{wrapfigure}
}




% Estas definiciones son para el comando \InsertFigBox
\newlength{\anchoFigura}
\newlength{\anchoFloat}
\addtolength{\fboxsep}{2\fboxsep}
%\renewcommand{\capfont}{\normalfont\normalcolor\sffamily\small}
%\renewcommand{\caplabelfont}{\normalfont\normalcolor\sffamily\bfseries\small}

% El comando \InsertFigBox nos permite insertar figuras en un marco
% Los parametros son:
% 1 --> Fichero de la imagen
% 2 --> Etiqueta (label) para referencias
% 3 --> Texto a pie de imagen
% 4 -> Porcentaje del ancho de página que ocupará la figura (de 0 a 1)
% 5 --> Opciones que queramos pasarle al \includegraphics
\newcommand{\InsertFigBox}[5]{%
  \setlength{\anchoFloat}{#4\textwidth}%
  \addtolength{\anchoFloat}{-4\fboxsep}%
  \setlength{\anchoFigura}{\anchoFloat}%
  \begin{figure}%
    \begin{center}%
      \Ovalbox{%
        \begin{minipage}{\anchoFloat}%
          \begin{center}%
            \includegraphics[width=\anchoFigura,#5]{#1}%
            \caption{#3}%
            \label{#2}%
          \end{center}%
        \end{minipage}
      }%
    \end{center}%
  \end{figure}%
}



%---------------------------------------------------------------
% Misc
%---------------------------------------------------------------

% predefined commands by Harish
\newcommand{\PdfPsText}[2]{
  \ifpdf
     #1
  \else
     #2
  \fi
}


%---------------------------------------------------------------
% Locales
%---------------------------------------------------------------


%%
%% Para quitar traducciones raras (Cuadros)
%% A de usarse cada vez que se seleccione el idioma
%%
\newcommand{\MejorarTraducciones}{%
       \renewcommand{\listtablename}{Índice de tablas}
       \renewcommand{\tablename}{Tabla}
       \renewcommand{\lstlistingname}{Lista}
}%



%---------------------------------------------------------------
% Source code
%---------------------------------------------------------------


%%
%% Para escribir extractos de codigo
%%
%% Las tabulaciones se substituyen por dos espacios
%\fvset{tabsize=2}
%% Creamos un nuevo environment de fancyvrb para los ejemplos enmarcados
%\DefineVerbatimEnvironment{VerbEj}{BVerbatim}{fontsize=\small,samepage=true,commandchars=\\\{\}}
%% Colo de fondo
%\definecolor{grisfondo}{gray}{0.9}
%% Environment para extractos de codigo
%\newenvironment{codigo}%
%{\VerbatimEnvironment\begin{Sbox}\begin{VerbEj}}%
%{\end{VerbEj}\end{Sbox}\setlength{\fboxsep}{8pt}\begin{center}\fcolorbox{black}{grisfondo}{\TheSbox}\end{center}}
%
%% Otro formato más bonito para código fuente
%\newcommand{\codigofuente}[3]{%
%  \lstinputlisting[language=#1,caption={#2}]{#3}%
%}





% My cousin *really* hates overindentations in descriptions
\setlist[description]{labelindent=5pt,leftmargin=0.5cm} % ,style=multiline


%: ----------------------------------------------------------------------
%:                  TITLE PAGE: name, degree,..
% ----------------------------------------------------------------------
% below is to generate the title page with crest and author name

% if output to PDF then put the following in PDF header
\ifpdf  
    \pdfinfo { /Title  (PhD)
               /Creator (TeX)
               /Producer (pdfTeX)
               /Author (Name surname)
               /CreationDate (D:YYYYMMDDhhmmss)  %format D:YYYYMMDDhhmmss
               /ModDate (D:YYYYMMDDhhmm)
               /Subject (xyz)
               /Keywords (keyword1, keyword2, keyword3) }
    \pdfcatalog { /PageMode (/UseOutlines)
                  /OpenAction (fitbh)  }
\fi





% ----------------------------------------------------------------------
% This section below defines front covert (external and internal)
% Shield logo
\crest{\includegraphics[width=2cm]{Deusto_Shield}}
% Full logo
%\crest{\includegraphics[width=6cm]{UDeusto}}
\makeatletter
% check http://tex.stackexchange.com/questions/27710/how-can-i-use-author-date-and-title-after-maketitle
% Title of the dissertation
\title{Semantic Tuples Spaces for Constrained Devices: A Web-compliant Vision} \let\Title\@title
\university{Universidad de Deusto}
\degree{Tesis doctoral presentada por }
\author{Aitor Gómez Goiri} \let\Author\@author
\collegeordept{dentro del Programa de Doctorado en Your programme name}
\textadvisor{Dirigida por }
\advisor{Dr. Diego López de Ipiña} \let\Advisor\@advisor
\textadvisortwo{y}
\advisortwo{Dr. Íñigo Goiri Presa} \let\Advisortwo\@advisortwo
\textsignaturecandidate{El doctorando}
\textsignatureadvisor{El director}
\cityofbirth{Bilbao}
%\degreedate{\monthname \ \the\year}
\degreedate{Month de \the\year}
\makeatother

\newcommand{\website}{\url{http://gomezgoiri.net}}


% ----------------------------------------------------------------------
       
% turn of those nasty overfull and underfull hboxes
\hbadness=10000
\hfuzz=50pt


%: --------------------------------------------------------------
%:                  FRONT MATTER: dedications, abstract,..
% --------------------------------------------------------------

\begin{document}

\selectlanguage{british}

% sets line spacing
\renewcommand\baselinestretch{1.2}
\baselineskip=18pt plus1pt

% Watermark
%\watermark{DRAFT	DRAFT	DRAFT	DRAFT	DRAFT	DRAFT	DRAFT	DRAFT	DRAFT}


%: ----------------------- generate cover page ------------------------

\maketitle  % command to print the title page with above variables



% Title back
% Thesis Titleback ---------------------------------------------------

\thispagestyle{empty}

\hfill

\vfill

\medskip


\noindent
\textit{
\Title
}




Author: \Author

Advisor: \Advisor

Advisor: \Advisortwo



\vfill

\vfill

\noindent
The following web-page address contains up to date information about this dissertation and related topics: \\
\website


\noindent
Text printed in Bilbao

% TODO final date
\noindent
First edition, 
% Moth and year
\monthname \ \the\year

\vspace{1cm}
\hrule
\bigskip

% \cleardoublepage command ends the current page and causes all figures and tables that have so far appeared in the input to be printed. In a two-sided printing style, it also makes the next page a right-hand (odd-numbered) page, producing a blank page if necessary. 
\cleardoublepage

%%: ----------------------- cover page back side ------------------------
%% Your research institution may require reviewer names, etc.
%% This cover back side is required by Dresden Med Fac; uncomment if needed.
%
%\newpage
%\vspace{10mm}
%1. Reviewer: Name
%
%\vspace{10mm}
%2. Reviewer: 
%
%\vspace{20mm}
%Day of the defense:
%
%\vspace{20mm}
%\hspace{70mm}Signature from head of PhD committee:
%
%
%\cleardoublepage

% ----------------------------------------------------------------------



%: ----------------------- abstract ------------------------

% Your institution may have specific regulations if you need an abstract and where it is to be placed in the document. The default here is just after title.


% The original template provides and abstractseparate environment, if your institution requires them to be separate. I think it's easier to print the abstract from the complete thesis by restricting printing to the relevant page.
% \begin{abstractseparate}
%   
% Thesis Abstract -----------------------------------------------------


%\begin{abstractslong}    %uncommenting this line, gives a different abstract heading


\begin{abstracts}        %this creates the heading for the abstract page
% Ubicomp, mobile e IoT
\acf{ubicomp} envisions environments where devices interact among them to seamlessly work on behalf of the humans.
% TODO TODO TODO In recent years,... in our days.
In recent years, the emergence of the \acf{iot} concept, which opts for connecting everyday objects to the Internet, and the Mobile Computing paradigm have contributed to strengthen UbiComp in our days.
For this reason, UbiComp environments are not necessarily populated by powerful computers.
On the contrary, mainly resource constrained devices (e.g., embedded and mobile devices) are the main actors of these environments.
% Retos de Ubicomp
Thus, it is important for the environment to deal with their heterogeneity, unreliability, and replaceability.


% Tuplespaces y semántica
To face heterogeneity, the Semantic Web has proposed several standards and models to provide a precise meaning to the terms so that they can be reused across applications boundaries.
Regarding unreliability and replaceability, space-based computing (or Tuple Spaces) promotes the uncoupled coordination of the devices.
Solutions based on semantic tuple spaces combine these three beneficial aspects, resulting from bridging the Semantic Web and Tuple Spaces domains, for UbiComp.

% Reto que supone adaptar el paradigma para cumplir los objetivos de la hipótesis
Most of these semantic tuple spaces consider embedded and mobile devices as mere clients of a space managed by more powerful devices.
Such delegation helps to reduce the workload of devices with computing and energy limitations.
However,  this delegation moves the data away from where it is physically generated.
This creates a conflict between providing updated data and generating unnecessary network traffic for unused information.
% TODO TODO TODO each other no!
Besides, this delegation makes constrained devices intrinsically dependent on each other when it might not always be necessary.
This dissertation explores how these constrained devices can act as fully fledged semantic knowledge providers to create a more decentralized space.

% Aspectos que se analizan en la tesis
In conclusion, this dissertation presents a novel adaptation of semantic tuple space which considers the energy and computational impact on the devices.
Specifically, this dissertation proposes the following contributions:
\begin{itemize}
  \item A space model which considers the principles which made the web flourish in the last decades, together with the uncoupling properties of the space-based computing.
  \item An energy-aware search mechanism for autonomous constrained devices.
  \item An alignment of two approaches to act on the physical environment, namely a space-based indirect actuation and a web-based direct actuation.
\end{itemize}

\end{abstracts}



\begin{resumen}        %this creates the heading for the abstract page
La computación ubicua (\acs{ubicomp}) vislumbra entornos donde los dispositivos interactuan entre sí para trabajar de forma imperceptible para los seres humanos.
En los últimos años, la emergencia del Internet de las Cosas (\acs{iot}), que aboga por conectar objetos cotidianos a Internet, y la computación móvil han contribuído a fortalecer \ac{ubicomp}.
Por ello, los entornos ubicuos no se componen necesariamente de computadoras potentes.
Por el contrario, los actores principales de dichos entornos suelen ser en su mayor parte dispositivos con recursos limitados (p.e. dispositivos móviles y embebidos).
Es por ello que es importante para estos entornos enfrentarse a la heterogeneidad, falta de fiabilidad y facilidad de reemplazo de estos dispositivos.


Para enfrentarse a la heterogeneidad, la web semántica propone diversos estándares y modelos para proveer un significado precido a los términos para que puedan reusarse más allá de las fronteras marcadas por las aplicaciones.
Relativo a la falta de fiabilidad y facilidad de reemplazo, la computación basada en espacios (o \acl{ts}) promueve una coordinación desacoplada de los dispositivos.
Las soluciones basadas en espacios de tuplas semánticos combinan estos tres aspectos beneficiosos para \ac{ubicomp} uniendo los dominios de la web semántica y los \acl{ts}.


Muchos de estos espacios de tuplas semánticos consideran a los dispositivos móviles y embebidos como meros clientes de un espacio gestionado por dispositivos más potentes.
Dicha delegación ayuda a reducir la carga de trabajo de los dispositivos con limitaciones computacionales y energéticas.
Sin embargo, esta delegación distancia los datos de donde son generados.
Esto crea un conflicto entre proveer datos actualizados y generar tráfico de red innecesario para información no usada.
Además, esta delegación hace a los dispositivos limitados intrinsicamente dependiente de otros cuando no siempre es necesario.
Esta tesis explora cómo esos dispositivos limitados pueden actuar como autenticos proveedores de conocimiento semántico para crear un espacio más descentralizado.


En conclusión, esta tesis presenta una adaptación novedora de los espacios de tuplas semánticos que considera el impacto computacional y energético en los dispositivos.
Especificamente, esta disertación propone las siguientes contribuciones:
\begin{itemize}
  \item Un modelo de espacio que considera los principios que han hecho florecer a la web en las últimas décadas, junto con las propiedades de desacoplamiento de la computación basada en espacios.
  \item Un mecanismo de búsqueda energéticamente eficiente para dispositivos limitados autónomos.
  \item Un alineamiento entre dos formas de actuar en un entorno físico: actuación indirecta basada en espacios y actuación directa basada en la web.
\end{itemize}

\end{resumen}



\begin{laburpena}        %this creates the heading for the abstract page
% Jarri zure laburpena hemen.


Vivamus ullamcorper erat in sapien dignissim pellentesque.


\end{laburpena}

%\end{abstractlongs}


% ---------------------------------------------------------------------- 

% \end{abstractseparate}


%: ----------------------- tie in front matter ------------------------

% The frontmatter text starts here
\frontmatter

% Thesis Dedictation ---------------------------------------------------

\begin{dedication} %this creates the heading for the dedication page

\textit{
Se lo dedico a todo doctorando que no vea la luz al final del tunel o se sienta desbordado por la frustración. Ánimo.
}

% TODO Dibujo de PhD comics

\emph{Sometimes, there's so much beauty in the world... I feel like I can't take it, like my heart is just going to cave in.}
\qauthor{Ricky Fitts, American Beauty}

\end{dedication}

% ----------------------------------------------------------------------


% Thesis Abstract -----------------------------------------------------


%\begin{abstractslong}    %uncommenting this line, gives a different abstract heading


\begin{abstracts}        %this creates the heading for the abstract page
% Ubicomp, mobile e IoT
\acf{ubicomp} envisions environments where devices interact among them to seamlessly work on behalf of the humans.
% TODO TODO TODO In recent years,... in our days.
In recent years, the emergence of the \acf{iot} concept, which opts for connecting everyday objects to the Internet, and the Mobile Computing paradigm have contributed to strengthen UbiComp in our days.
For this reason, UbiComp environments are not necessarily populated by powerful computers.
On the contrary, mainly resource constrained devices (e.g., embedded and mobile devices) are the main actors of these environments.
% Retos de Ubicomp
Thus, it is important for the environment to deal with their heterogeneity, unreliability, and replaceability.


% Tuplespaces y semántica
To face heterogeneity, the Semantic Web has proposed several standards and models to provide a precise meaning to the terms so that they can be reused across applications boundaries.
Regarding unreliability and replaceability, space-based computing (or Tuple Spaces) promotes the uncoupled coordination of the devices.
Solutions based on semantic tuple spaces combine these three beneficial aspects, resulting from bridging the Semantic Web and Tuple Spaces domains, for UbiComp.

% Reto que supone adaptar el paradigma para cumplir los objetivos de la hipótesis
Most of these semantic tuple spaces consider embedded and mobile devices as mere clients of a space managed by more powerful devices.
Such delegation helps to reduce the workload of devices with computing and energy limitations.
However,  this delegation moves the data away from where it is physically generated.
This creates a conflict between providing updated data and generating unnecessary network traffic for unused information.
% TODO TODO TODO each other no!
Besides, this delegation makes constrained devices intrinsically dependent on each other when it might not always be necessary.
This dissertation explores how these constrained devices can act as fully fledged semantic knowledge providers to create a more decentralized space.

% Aspectos que se analizan en la tesis
In conclusion, this dissertation presents a novel adaptation of semantic tuple space which considers the energy and computational impact on the devices.
Specifically, this dissertation proposes the following contributions:
\begin{itemize}
  \item A space model which considers the principles which made the web flourish in the last decades, together with the uncoupling properties of the space-based computing.
  \item An energy-aware search mechanism for autonomous constrained devices.
  \item An alignment of two approaches to act on the physical environment, namely a space-based indirect actuation and a web-based direct actuation.
\end{itemize}

\end{abstracts}



\begin{resumen}        %this creates the heading for the abstract page
La computación ubicua (\acs{ubicomp}) vislumbra entornos donde los dispositivos interactuan entre sí para trabajar de forma imperceptible para los seres humanos.
En los últimos años, la emergencia del Internet de las Cosas (\acs{iot}), que aboga por conectar objetos cotidianos a Internet, y la computación móvil han contribuído a fortalecer \ac{ubicomp}.
Por ello, los entornos ubicuos no se componen necesariamente de computadoras potentes.
Por el contrario, los actores principales de dichos entornos suelen ser en su mayor parte dispositivos con recursos limitados (p.e. dispositivos móviles y embebidos).
Es por ello que es importante para estos entornos enfrentarse a la heterogeneidad, falta de fiabilidad y facilidad de reemplazo de estos dispositivos.


Para enfrentarse a la heterogeneidad, la web semántica propone diversos estándares y modelos para proveer un significado precido a los términos para que puedan reusarse más allá de las fronteras marcadas por las aplicaciones.
Relativo a la falta de fiabilidad y facilidad de reemplazo, la computación basada en espacios (o \acl{ts}) promueve una coordinación desacoplada de los dispositivos.
Las soluciones basadas en espacios de tuplas semánticos combinan estos tres aspectos beneficiosos para \ac{ubicomp} uniendo los dominios de la web semántica y los \acl{ts}.


Muchos de estos espacios de tuplas semánticos consideran a los dispositivos móviles y embebidos como meros clientes de un espacio gestionado por dispositivos más potentes.
Dicha delegación ayuda a reducir la carga de trabajo de los dispositivos con limitaciones computacionales y energéticas.
Sin embargo, esta delegación distancia los datos de donde son generados.
Esto crea un conflicto entre proveer datos actualizados y generar tráfico de red innecesario para información no usada.
Además, esta delegación hace a los dispositivos limitados intrinsicamente dependiente de otros cuando no siempre es necesario.
Esta tesis explora cómo esos dispositivos limitados pueden actuar como autenticos proveedores de conocimiento semántico para crear un espacio más descentralizado.


En conclusión, esta tesis presenta una adaptación novedora de los espacios de tuplas semánticos que considera el impacto computacional y energético en los dispositivos.
Especificamente, esta disertación propone las siguientes contribuciones:
\begin{itemize}
  \item Un modelo de espacio que considera los principios que han hecho florecer a la web en las últimas décadas, junto con las propiedades de desacoplamiento de la computación basada en espacios.
  \item Un mecanismo de búsqueda energéticamente eficiente para dispositivos limitados autónomos.
  \item Un alineamiento entre dos formas de actuar en un entorno físico: actuación indirecta basada en espacios y actuación directa basada en la web.
\end{itemize}

\end{resumen}



\begin{laburpena}        %this creates the heading for the abstract page
% Jarri zure laburpena hemen.


Vivamus ullamcorper erat in sapien dignissim pellentesque.


\end{laburpena}

%\end{abstractlongs}


% ---------------------------------------------------------------------- 


% Thesis Acknowledgements ------------------------------------------------


% Opening of the acknowledgements

%Sort version
%this creates the heading for the acknowlegments

\begin{acknowledgements}      
%Long version
%uncommenting this line, gives a different acknowledgements heading
%\begin{acknowledgementslong} 

\begin{flushright}
\emph{Sometimes, there's so much beauty in the world - I feel like I can't take it, like my heart is just going to cave in.}
\qauthor{Ricky Fitts (American Beauty, 1999).}
\end{flushright}

\bigskip

Si debo a alguien la posibilidad de haber concluido esta aventura es sin duda a aita y ama.
Esta posibilidad va más allá del mundano y sin embargo tan necesario hecho de haberme proporcionado una cómoda vida en la que siempre pude elegir mi camino.
Por vosotros soy como soy y soy quien soy.
Podría haber salido un hijo un poco más hablador y menos pedante, pero no será porque no os esforzasteis.
Echando la vista atrás no puedo sino estar orgullosísimo de la educación recibida.
Por ello, vaya para vosotros mi primer agradecimiento.


Gracias Diego por haberme dado esta enorme oportunidad.
Gracias por el tiempo y el espacio concedido.
Gracias Íñigo por soportarme en momentos de máxima incertidumbre y pesimismo.
Gracias por servirme de guía.


Gracias a todos mis compañeros y ex-compañeros.
Gracias a mis compañeros de generación por compartir sufrimientos: Edu, Iván y Eguiluz.
Gracias a esa infatigable e incansable máquina con piernas de la que tanto se aprende que es Pablo Orduña.
Gracias a los ancianos del lugar: Aitor y Unai.
Gracias a la gente del Garilab por aceptarme durante mi exilio.
En definitiva, gracias a tanta gente que ha pasado por la vida \emph{esmarlabense} estos últimos años: Don Jaime, Emaldi, Laiseca, el innombrable, Lázaro y un largo etc.
Afortunadamente, ese ``\emph{largo etc}'' se puede consultar en la página web del grupo haciendo uso de la consulta mostrada al final de esta sección o siguiendo el siguiente \emph{QRCode}.
% Poner el SPARQL para seleccionarlos!



\begin{figure}[h]
    \centering
    \includegraphics[width=0.3\textwidth]{qrcode_acks}
\end{figure}


Thanks to all the inspiring people I've known during this journey.
Thanks to Geoff Coulson, Barry Porter and the rest of the people I met at Lancaster University.
Thanks to all the ETH's \emph{Distributed Systems group} members, specially to Simon Mayer, my cicerone in Zurich.
Thanks to rest of you, no matter if I've met you in person or not.


Gracias a todos mis compañeros de aventuras estos años.
A los ocasionales y a los fieles.
En viajes por el mundo, canchas de baloncesto, senderos, \emph{eventos del esmarla},...
Cada segundo en que mantuve la cabeza fuera de la tesis fue un respiro enorme.


Gracias al resto de mi familia, la sanguínea y la \emph{circunstancial}.
Especialmente a mi hermano Mikel y a mi abuela Teo.


\begin{flushright}
\textit{Eskerrik asko,}

Aitor Gómez Goiri

% Month and year
\monthname \ \the\year

% Signature figure

%\begin{figure}[htbp!]
%\end{figure}
%\includegraphics{signature}%

\end{flushright}

\begin{listing}[nolol]
  \scriptsize
  \expandafter\def\csname PY@tok@err\endcsname{}
\begin{Verbatim}[commandchars=\\\{\}]
\PY{k}{PREFIX }\PY{n+nv}{swrcfe:}\PY{n+nn}{ \PYZlt{}http://www.morelab.deusto.es/ontologies/swrcfe.owl\PYZsh{}\PYZgt{}}

\PY{k}{SELECT }\PY{err}{C}\PY{err}{O}\PY{err}{U}\PY{err}{N}\PY{err}{T}\PY{err}{(}\PY{err}{D}\PY{err}{I}\PY{err}{S}\PY{err}{T}\PY{err}{I}\PY{err}{N}\PY{err}{C}\PY{err}{T}\PY{err}{ }\PY{err}{?}\PY{err}{f}\PY{err}{a}\PY{err}{n}\PY{err}{a}\PY{err}{m}\PY{err}{e}\PY{err}{)} \PY{k}{WHERE}\PY{p}{\PYZob{}}
\PY{n+nv}{  ?job}\PY{o}{ swrcfe:doneBy}
\PY{o}{     }\PY{n+na}{\PYZlt{}http://www.morelab.deusto.es/labman/resource/people/aitor\PYZhy{}gomez\PYZhy{}goiri\PYZgt{}} \PY{p}{.}
  \PY{p}{\PYZob{}}
    \PY{n+nv}{?job2}\PY{o}{ swrcfe:doneBy }\PY{n+nv}{?formerworkmate }\PY{p}{.}
    \PY{n+nv}{?formerworkmate}\PY{o}{ foaf:familyName }\PY{n+nv}{?faname }\PY{p}{.}
    \PY{n+nv}{?formerworkmate}\PY{o}{ foaf:firstName }\PY{n+nv}{?finame }\PY{p}{.}
    \PY{n+nv}{?job}\PY{o}{ swrcfe:jobStartDate }\PY{n+nv}{?jobstartdate }\PY{p}{.}

    \PY{k}{OPTIONAL} \PY{p}{\PYZob{}}\PY{n+nv}{ ?job2}\PY{o}{ swrcfe:jobEndDate  }\PY{n+nv}{?jobenddate }\PY{err}{\PYZcb{}}
\PY{n+na}{    FILTER}\PY{p}{(}\PY{n+nv}{?jobstartdate }\PY{err}{\PYZlt{}}\PY{err}{ }\PY{n+nv}{?jobenddate}\PY{p}{)}
  \PY{p}{\PYZcb{}}\PY{err}{ }\PY{err}{U}\PY{err}{N}\PY{err}{I}\PY{err}{O}\PY{err}{N} \PY{p}{\PYZob{}}
\PY{n+nv}{    ?job2}\PY{o}{ swrcfe:doneBy }\PY{n+nv}{?workmate }\PY{p}{.}
    \PY{n+nv}{?workmate}\PY{o}{ foaf:familyName }\PY{n+nv}{?faname }\PY{p}{.}
    \PY{n+nv}{?workmate}\PY{o}{ foaf:firstName }\PY{n+nv}{?finame }\PY{p}{.}

    \PY{k}{OPTIONAL} \PY{p}{\PYZob{}}\PY{n+nv}{ ?job2}\PY{o}{ swrcfe:jobEndDate }\PY{n+nv}{?jobenddate }\PY{err}{\PYZcb{}}
\PY{n+na}{    FILTER}\PY{p}{(}\PY{err}{!}\PY{n+na}{bound}\PY{p}{(}\PY{n+nv}{?jobenddate}\PY{p}{)}\PY{p}{)}
  \PY{p}{\PYZcb{}}
\PY{err}{ }\PY{err}{ }\PY{err}{F}\PY{err}{I}\PY{err}{L}\PY{err}{T}\PY{err}{E}\PY{err}{R}\PY{err}{(}\PY{err}{?}\PY{err}{j}\PY{err}{o}\PY{err}{b}\PY{err}{ }\PY{err}{!}\PY{err}{=}\PY{err}{ }\PY{err}{?}\PY{err}{j}\PY{err}{o}\PY{err}{b}\PY{err}{2}\PY{err}{)}
\PY{p}{\PYZcb{}}\PY{err}{ }
\end{Verbatim}

  %\caption{Consulta SPARQL para obtener la gente relacionada con el grupo MORElab a la que tengo que estar agradecido.}
  \label{lst:workmates_ack}
\end{listing}



%Closing of the acknowledgements
%Sort version
\end{acknowledgements}
% Long version
%\end{acknowledgementslong}

% ------------------------------------------------------------------------




% As abstract contains various languages we set the main language again
\selectlanguage{british}


%: ----------------------- contents ------------------------

\setcounter{secnumdepth}{5} % organisational level that receives a numbers
\setcounter{tocdepth}{2}    % print table of contents for level 3
% "Note that if you have a list of figures/tables, the tocdepth must be at least 2 for them to show up"

\hypersetup{ % requires hyperref package
  bookmarksdepth = 3, % for the PDF's table of contents!
}


%%You can also add extra lines to the ToC or to force extra unnumbered section headings to be included. For example, if you wanted to add an entry called Preface, and you didn't want the Preface to be numbered, you'd use these commands:
%\ subsection*{Preface}
%\addcontentsline{toc}{subsection}{Preface} 

\tableofcontents            % print the table of contents
% levels are: 0 - chapter, 1 - section, 2 - subsection, 3 - subsubsection

%: ----------------------- list of figures/tables ------------------------

\listoffigures	% print list of figures
\listoftables  % print list of tables


%: ----------------------- glossary ------------------------

% Usando ACRONYMS, puedo definir palabras rollo Tuple Spaces para ser uniforme a lo largo del texto.
% Así evito acabar llamandolas tspaces, tuplespaces, o tuple spaces.
% Comprobar sobre texto escrito que uso variantes de \ac siempre
% Faltan por comprobar los capítulos 4, 5 y 6

%% acronyms
%% http://staff.science.uva.nl/~polko/HOWTO/LATEX/acronym.html
%\acrodef{label}[acronym]{written out form}

\acrodef{ts}[TS]{Tuple Space}
\acrodef{tsc}[TSC]{Triple Space Computing}
\acrodef{ami}[AmI]{Ambient Intelligence}
\acrodef{ubicomp}[UbiComp]{ubiquitous computing}
\acrodef{wot}[WoT]{Web of Things}
\acrodef{iot}[IoT]{Internet of Things}
\acrodef{sw}[SW]{Semantic Web}


% si decido gustarme: http://www.latex-community.org/know-how/263-glossaries-nomenclature-lists-of-symbols-and-acronyms
% Tie in external source file for definitions: /0_frontmatter/glossary.tex
% Glossary entries can also be defined in the main text. See glossary.tex
% this file is called up by thesis.tex
% content in this file will be fed into the main document

% Glossary entries are defined with the command \nomenclature{1}{2}
% 1 = Entry name, e.g. abbreviation; 2 = Explanation
% You can place all explanations in this separate file or declare them in the middle of the text. Either way they will be collected in the glossary.

% required to print nomenclature name to page header
\markboth{\MakeUppercase{\nomname}}{\MakeUppercase{\nomname}}

% ----------------------- contents from here ------------------------
%

%
%
%% acronyms

\nomenclature{AmI}{Ambient Intelligence}
\nomenclature{IP}{Internet Protocol}
\nomenclature{WoT}{Web of Things}
\nomenclature{TS}{Tuple Space}
\nomenclature{TSC}{Triple Space Computing}
\nomenclature{UbiComp}{Ubiquitous Computing}




%\begin{multicols}{2} % \begin{multicols}{#columns}[header text][space]
%\begin{footnotesize} % scriptsize(7) < footnotesize(8) < small (9) < normal (10)

%\printnomenclature[1.5cm] % [] = distance between entry and description

% I don't know why it does not work.
% TODO fix or remove (acronyms are already listed)
\printnomenclature % [] = distance between entry and description

\label{sec:glossary} % target name for links to glossary

%\end{footnotesize}
%\end{multicols}




%: --------------------------------------------------------------
%:                  MAIN DOCUMENT SECTION
% --------------------------------------------------------------

% the main text starts here with the introduction, 1st chapter,...
\mainmatter

%\renewcommand{\chaptername}{} % uncomment to print only "1" not "Chapter 1"
\pagestyle{fancy}


%: ----------------------- subdocuments ------------------------

% Parts of the thesis are included below. Rename the files as required.
% But take care that the paths match. You can also change the order of appearance by moving the include commands.

\newcommand{\pathchapone}{1_introduction}
\newcommand{\pathchaptwo}{2_background}
\newcommand{\pathchapthree}{3_state_of_the_art}
\newcommand{\pathchapfour}{4_tsc}
\newcommand{\pathchapfive}{5_search}
\newcommand{\pathchapsix}{6_actuate}
\newcommand{\pathchapseven}{7_conclusion}


%: ----------------------- introduction ------------------------

% this file is called up by thesis.tex
% content in this file will be fed into the main document

%: ----------------------- introduction file header -----------------------
\begin{savequote}[50mm]
I tell you what, I speak slow so those of you with a Phd can understand.
\qauthor{Doug \textit{(Déjà Vu)}}
\end{savequote}


\chapter{Conclusions}
\label{cha:conclusions}
\newcommand{\pathchapseven}{7_conclusion}


% the code below specifies where the figures are stored
\ifpdf
    \graphicspath{{\pathchapseven/figures/PNG/}{\pathchapseven/figures/PDF/}{\pathchapseven/figures/}}
\else
    \graphicspath{{\pathchapseven/figures/EPS/}{\pathchapseven/figures/}}
\fi


%------------------------------------------------------------------------- 

% Presentar de qué ha ido la tesis
% 3 ptos con los que quiero que el lector se quede: searching, acting y compatibility with WoT

This dissertation explores the design of a completely distributed \ac{tsc} middleware.
In this middleware, each node is responsible of its own information, but together with other nodes they create a virtual space. % virtual, concept
Writing and reading from this space, the nodes communicate in a much more decoupled indirect style.

The distribution of the information avoids the dependency of the participants on other devices, but arises other challenges.
We have tackled two main challenges: how to search in a decentralized way and how to act on the environment through the space.
Besides, after an analysis we acknowledged the importance and benefits that the upcoming \ac{wot} has brought to the \acl{ubicomp}.
This has lead us to have make our final solution compatible with the \ac{wot} and other REST services.

To enhance the search process between devices limiting their dependencies in others we designed an architecture.
This architecture dynamically chooses an intermediary according to its capabilities to support the most limited devices.
The final result is that less requests are required in overall, specifically reducing the load on resource constrained devices.

% TODO refinar cuando esté hecha esa parte
We also presented two different acting mechanisms: X and Y.
Comparing them we reach to the conclusion that X proposes a less decoupled design, which in turn depends on a subscription mechanism.
This subscription mechanism is particularly challenging on dynamic environments.
The Y solution is more straightforward and more based in REST services consumption.
In contrast, it depends on reasoners and on some additional information about the services offered by the rest of the nodes.

Besides, we have guided our design to build the middleware over a REST interface.
Using that interface, each node exposes the semantic resources, i.e. RDF triples grouped in RDF graphs, it manages.
This enables the consumption of these contents by other semantic \ac{wot} solutions or REST services, even if they are not part of our middleware.

\section{Contributions}

% De una charla con Buján recuerdo que dijo que para O. Corcho...
%   + Contribuciones científicas: lo que nadie había hecho antes (innovación)
%   + Contribuciones técnicas: implementar o aplicar en otro campo un aspecto innovado por otro

% Contributions (por si necesito ideas)
% \item Analysis of strengths and weaknesses of both \ac{wot} and \ac{tsc} and comparison of both approaches.
% \item Design and implementation of a RESTful interface which makes possible to build a \ac{tsc} on top of it. % TODO revisar y poner todos los peros que haga falta para no traicionar la idea de Fielding
%       This resource-oriented interface enables the access to the the semantic knowledge provided by each device.
%       By doing so, we keep our middleware compliant with the \ac{wot} idea. % o integrable con otros servicios, o...
%       % es como explicar una dualidad: puedes usar nuestro middleware usando las primitivas de TSC, pero también puedes usarlo simplemente porque te viene bien para WoT semántico
% \item A search-aware architecture which promotes the end-to-end communication between devices. % and balances their load by administrative tasks according to their capacities.
%       This architecture tries to minimize the requests devices have to handle by enhancing their search mechanism.
%       The extra-tasks introduced by this enhancement are performed by nodes chosen according to their capacities.
% \item Design and comparison of two approaches which enable to act on the environment by writing in the fully distributed space.
% \item Exploration of the use of semantic content in several current embedded platforms.
%       We have tested several platforms, semantic frameworks and reasoners to verify the feasibility of the ideas presented in the dissertation.
% \item Design and implementation of most of the ideas presented in a middleware publicly available called \emph{Otsopack}. % al primo no le gustará, pero insisto
%       This middleware has been used in several projects for a range of scenarios.


\section{Relevant Publications}

% International JCR Journals
% International conferences
% Workshops % o si queda muy corto, incluirlo directamente en conclusions
\section{Future Work}

% Limitations => Future Work
\section{Final Remarks}



% ----------------------------------------------------------------------

%: ----------------------- related work ------------------------

% this file is called up by thesis.tex
% content in this file will be fed into the main document

%: ----------------------- introduction file header -----------------------
\begin{savequote}[50mm]
I tell you what, I speak slow so those of you with a Phd can understand.
\qauthor{Doug \textit{(Déjà Vu)}}
\end{savequote}


\chapter{Conclusions}
\label{cha:conclusions}
\newcommand{\pathchapseven}{7_conclusion}


% the code below specifies where the figures are stored
\ifpdf
    \graphicspath{{\pathchapseven/figures/PNG/}{\pathchapseven/figures/PDF/}{\pathchapseven/figures/}}
\else
    \graphicspath{{\pathchapseven/figures/EPS/}{\pathchapseven/figures/}}
\fi


%------------------------------------------------------------------------- 

% Presentar de qué ha ido la tesis
% 3 ptos con los que quiero que el lector se quede: searching, acting y compatibility with WoT

This dissertation explores the design of a completely distributed \ac{tsc} middleware.
In this middleware, each node is responsible of its own information, but together with other nodes they create a virtual space. % virtual, concept
Writing and reading from this space, the nodes communicate in a much more decoupled indirect style.

The distribution of the information avoids the dependency of the participants on other devices, but arises other challenges.
We have tackled two main challenges: how to search in a decentralized way and how to act on the environment through the space.
Besides, after an analysis we acknowledged the importance and benefits that the upcoming \ac{wot} has brought to the \acl{ubicomp}.
This has lead us to have make our final solution compatible with the \ac{wot} and other REST services.

To enhance the search process between devices limiting their dependencies in others we designed an architecture.
This architecture dynamically chooses an intermediary according to its capabilities to support the most limited devices.
The final result is that less requests are required in overall, specifically reducing the load on resource constrained devices.

% TODO refinar cuando esté hecha esa parte
We also presented two different acting mechanisms: X and Y.
Comparing them we reach to the conclusion that X proposes a less decoupled design, which in turn depends on a subscription mechanism.
This subscription mechanism is particularly challenging on dynamic environments.
The Y solution is more straightforward and more based in REST services consumption.
In contrast, it depends on reasoners and on some additional information about the services offered by the rest of the nodes.

Besides, we have guided our design to build the middleware over a REST interface.
Using that interface, each node exposes the semantic resources, i.e. RDF triples grouped in RDF graphs, it manages.
This enables the consumption of these contents by other semantic \ac{wot} solutions or REST services, even if they are not part of our middleware.

\section{Contributions}

% De una charla con Buján recuerdo que dijo que para O. Corcho...
%   + Contribuciones científicas: lo que nadie había hecho antes (innovación)
%   + Contribuciones técnicas: implementar o aplicar en otro campo un aspecto innovado por otro

% Contributions (por si necesito ideas)
% \item Analysis of strengths and weaknesses of both \ac{wot} and \ac{tsc} and comparison of both approaches.
% \item Design and implementation of a RESTful interface which makes possible to build a \ac{tsc} on top of it. % TODO revisar y poner todos los peros que haga falta para no traicionar la idea de Fielding
%       This resource-oriented interface enables the access to the the semantic knowledge provided by each device.
%       By doing so, we keep our middleware compliant with the \ac{wot} idea. % o integrable con otros servicios, o...
%       % es como explicar una dualidad: puedes usar nuestro middleware usando las primitivas de TSC, pero también puedes usarlo simplemente porque te viene bien para WoT semántico
% \item A search-aware architecture which promotes the end-to-end communication between devices. % and balances their load by administrative tasks according to their capacities.
%       This architecture tries to minimize the requests devices have to handle by enhancing their search mechanism.
%       The extra-tasks introduced by this enhancement are performed by nodes chosen according to their capacities.
% \item Design and comparison of two approaches which enable to act on the environment by writing in the fully distributed space.
% \item Exploration of the use of semantic content in several current embedded platforms.
%       We have tested several platforms, semantic frameworks and reasoners to verify the feasibility of the ideas presented in the dissertation.
% \item Design and implementation of most of the ideas presented in a middleware publicly available called \emph{Otsopack}. % al primo no le gustará, pero insisto
%       This middleware has been used in several projects for a range of scenarios.


\section{Relevant Publications}

% International JCR Journals
% International conferences
% Workshops % o si queda muy corto, incluirlo directamente en conclusions
\section{Future Work}

% Limitations => Future Work
\section{Final Remarks}



% ----------------------------------------------------------------------

%: ----------------------- related work ------------------------

% this file is called up by thesis.tex
% content in this file will be fed into the main document

%: ----------------------- introduction file header -----------------------
\begin{savequote}[50mm]
I tell you what, I speak slow so those of you with a Phd can understand.
\qauthor{Doug \textit{(Déjà Vu)}}
\end{savequote}


\chapter{Conclusions}
\label{cha:conclusions}
\newcommand{\pathchapseven}{7_conclusion}


% the code below specifies where the figures are stored
\ifpdf
    \graphicspath{{\pathchapseven/figures/PNG/}{\pathchapseven/figures/PDF/}{\pathchapseven/figures/}}
\else
    \graphicspath{{\pathchapseven/figures/EPS/}{\pathchapseven/figures/}}
\fi


%------------------------------------------------------------------------- 

% Presentar de qué ha ido la tesis
% 3 ptos con los que quiero que el lector se quede: searching, acting y compatibility with WoT

This dissertation explores the design of a completely distributed \ac{tsc} middleware.
In this middleware, each node is responsible of its own information, but together with other nodes they create a virtual space. % virtual, concept
Writing and reading from this space, the nodes communicate in a much more decoupled indirect style.

The distribution of the information avoids the dependency of the participants on other devices, but arises other challenges.
We have tackled two main challenges: how to search in a decentralized way and how to act on the environment through the space.
Besides, after an analysis we acknowledged the importance and benefits that the upcoming \ac{wot} has brought to the \acl{ubicomp}.
This has lead us to have make our final solution compatible with the \ac{wot} and other REST services.

To enhance the search process between devices limiting their dependencies in others we designed an architecture.
This architecture dynamically chooses an intermediary according to its capabilities to support the most limited devices.
The final result is that less requests are required in overall, specifically reducing the load on resource constrained devices.

% TODO refinar cuando esté hecha esa parte
We also presented two different acting mechanisms: X and Y.
Comparing them we reach to the conclusion that X proposes a less decoupled design, which in turn depends on a subscription mechanism.
This subscription mechanism is particularly challenging on dynamic environments.
The Y solution is more straightforward and more based in REST services consumption.
In contrast, it depends on reasoners and on some additional information about the services offered by the rest of the nodes.

Besides, we have guided our design to build the middleware over a REST interface.
Using that interface, each node exposes the semantic resources, i.e. RDF triples grouped in RDF graphs, it manages.
This enables the consumption of these contents by other semantic \ac{wot} solutions or REST services, even if they are not part of our middleware.

\section{Contributions}

% De una charla con Buján recuerdo que dijo que para O. Corcho...
%   + Contribuciones científicas: lo que nadie había hecho antes (innovación)
%   + Contribuciones técnicas: implementar o aplicar en otro campo un aspecto innovado por otro

% Contributions (por si necesito ideas)
% \item Analysis of strengths and weaknesses of both \ac{wot} and \ac{tsc} and comparison of both approaches.
% \item Design and implementation of a RESTful interface which makes possible to build a \ac{tsc} on top of it. % TODO revisar y poner todos los peros que haga falta para no traicionar la idea de Fielding
%       This resource-oriented interface enables the access to the the semantic knowledge provided by each device.
%       By doing so, we keep our middleware compliant with the \ac{wot} idea. % o integrable con otros servicios, o...
%       % es como explicar una dualidad: puedes usar nuestro middleware usando las primitivas de TSC, pero también puedes usarlo simplemente porque te viene bien para WoT semántico
% \item A search-aware architecture which promotes the end-to-end communication between devices. % and balances their load by administrative tasks according to their capacities.
%       This architecture tries to minimize the requests devices have to handle by enhancing their search mechanism.
%       The extra-tasks introduced by this enhancement are performed by nodes chosen according to their capacities.
% \item Design and comparison of two approaches which enable to act on the environment by writing in the fully distributed space.
% \item Exploration of the use of semantic content in several current embedded platforms.
%       We have tested several platforms, semantic frameworks and reasoners to verify the feasibility of the ideas presented in the dissertation.
% \item Design and implementation of most of the ideas presented in a middleware publicly available called \emph{Otsopack}. % al primo no le gustará, pero insisto
%       This middleware has been used in several projects for a range of scenarios.


\section{Relevant Publications}

% International JCR Journals
% International conferences
% Workshops % o si queda muy corto, incluirlo directamente en conclusions
\section{Future Work}

% Limitations => Future Work
\section{Final Remarks}



% ----------------------------------------------------------------------

%: ----------------------- TSC for resource constrained devices ------------------------

% this file is called up by thesis.tex
% content in this file will be fed into the main document

%: ----------------------- introduction file header -----------------------
\begin{savequote}[50mm]
I tell you what, I speak slow so those of you with a Phd can understand.
\qauthor{Doug \textit{(Déjà Vu)}}
\end{savequote}


\chapter{Conclusions}
\label{cha:conclusions}
\newcommand{\pathchapseven}{7_conclusion}


% the code below specifies where the figures are stored
\ifpdf
    \graphicspath{{\pathchapseven/figures/PNG/}{\pathchapseven/figures/PDF/}{\pathchapseven/figures/}}
\else
    \graphicspath{{\pathchapseven/figures/EPS/}{\pathchapseven/figures/}}
\fi


%------------------------------------------------------------------------- 

% Presentar de qué ha ido la tesis
% 3 ptos con los que quiero que el lector se quede: searching, acting y compatibility with WoT

This dissertation explores the design of a completely distributed \ac{tsc} middleware.
In this middleware, each node is responsible of its own information, but together with other nodes they create a virtual space. % virtual, concept
Writing and reading from this space, the nodes communicate in a much more decoupled indirect style.

The distribution of the information avoids the dependency of the participants on other devices, but arises other challenges.
We have tackled two main challenges: how to search in a decentralized way and how to act on the environment through the space.
Besides, after an analysis we acknowledged the importance and benefits that the upcoming \ac{wot} has brought to the \acl{ubicomp}.
This has lead us to have make our final solution compatible with the \ac{wot} and other REST services.

To enhance the search process between devices limiting their dependencies in others we designed an architecture.
This architecture dynamically chooses an intermediary according to its capabilities to support the most limited devices.
The final result is that less requests are required in overall, specifically reducing the load on resource constrained devices.

% TODO refinar cuando esté hecha esa parte
We also presented two different acting mechanisms: X and Y.
Comparing them we reach to the conclusion that X proposes a less decoupled design, which in turn depends on a subscription mechanism.
This subscription mechanism is particularly challenging on dynamic environments.
The Y solution is more straightforward and more based in REST services consumption.
In contrast, it depends on reasoners and on some additional information about the services offered by the rest of the nodes.

Besides, we have guided our design to build the middleware over a REST interface.
Using that interface, each node exposes the semantic resources, i.e. RDF triples grouped in RDF graphs, it manages.
This enables the consumption of these contents by other semantic \ac{wot} solutions or REST services, even if they are not part of our middleware.

\section{Contributions}

% De una charla con Buján recuerdo que dijo que para O. Corcho...
%   + Contribuciones científicas: lo que nadie había hecho antes (innovación)
%   + Contribuciones técnicas: implementar o aplicar en otro campo un aspecto innovado por otro

% Contributions (por si necesito ideas)
% \item Analysis of strengths and weaknesses of both \ac{wot} and \ac{tsc} and comparison of both approaches.
% \item Design and implementation of a RESTful interface which makes possible to build a \ac{tsc} on top of it. % TODO revisar y poner todos los peros que haga falta para no traicionar la idea de Fielding
%       This resource-oriented interface enables the access to the the semantic knowledge provided by each device.
%       By doing so, we keep our middleware compliant with the \ac{wot} idea. % o integrable con otros servicios, o...
%       % es como explicar una dualidad: puedes usar nuestro middleware usando las primitivas de TSC, pero también puedes usarlo simplemente porque te viene bien para WoT semántico
% \item A search-aware architecture which promotes the end-to-end communication between devices. % and balances their load by administrative tasks according to their capacities.
%       This architecture tries to minimize the requests devices have to handle by enhancing their search mechanism.
%       The extra-tasks introduced by this enhancement are performed by nodes chosen according to their capacities.
% \item Design and comparison of two approaches which enable to act on the environment by writing in the fully distributed space.
% \item Exploration of the use of semantic content in several current embedded platforms.
%       We have tested several platforms, semantic frameworks and reasoners to verify the feasibility of the ideas presented in the dissertation.
% \item Design and implementation of most of the ideas presented in a middleware publicly available called \emph{Otsopack}. % al primo no le gustará, pero insisto
%       This middleware has been used in several projects for a range of scenarios.


\section{Relevant Publications}

% International JCR Journals
% International conferences
% Workshops % o si queda muy corto, incluirlo directamente en conclusions
\section{Future Work}

% Limitations => Future Work
\section{Final Remarks}



% ----------------------------------------------------------------------

%: ----------------------- On how to Search in TSC ------------------------

% this file is called up by thesis.tex
% content in this file will be fed into the main document

%: ----------------------- introduction file header -----------------------
\begin{savequote}[50mm]
I tell you what, I speak slow so those of you with a Phd can understand.
\qauthor{Doug \textit{(Déjà Vu)}}
\end{savequote}


\chapter{Conclusions}
\label{cha:conclusions}
\newcommand{\pathchapseven}{7_conclusion}


% the code below specifies where the figures are stored
\ifpdf
    \graphicspath{{\pathchapseven/figures/PNG/}{\pathchapseven/figures/PDF/}{\pathchapseven/figures/}}
\else
    \graphicspath{{\pathchapseven/figures/EPS/}{\pathchapseven/figures/}}
\fi


%------------------------------------------------------------------------- 

% Presentar de qué ha ido la tesis
% 3 ptos con los que quiero que el lector se quede: searching, acting y compatibility with WoT

This dissertation explores the design of a completely distributed \ac{tsc} middleware.
In this middleware, each node is responsible of its own information, but together with other nodes they create a virtual space. % virtual, concept
Writing and reading from this space, the nodes communicate in a much more decoupled indirect style.

The distribution of the information avoids the dependency of the participants on other devices, but arises other challenges.
We have tackled two main challenges: how to search in a decentralized way and how to act on the environment through the space.
Besides, after an analysis we acknowledged the importance and benefits that the upcoming \ac{wot} has brought to the \acl{ubicomp}.
This has lead us to have make our final solution compatible with the \ac{wot} and other REST services.

To enhance the search process between devices limiting their dependencies in others we designed an architecture.
This architecture dynamically chooses an intermediary according to its capabilities to support the most limited devices.
The final result is that less requests are required in overall, specifically reducing the load on resource constrained devices.

% TODO refinar cuando esté hecha esa parte
We also presented two different acting mechanisms: X and Y.
Comparing them we reach to the conclusion that X proposes a less decoupled design, which in turn depends on a subscription mechanism.
This subscription mechanism is particularly challenging on dynamic environments.
The Y solution is more straightforward and more based in REST services consumption.
In contrast, it depends on reasoners and on some additional information about the services offered by the rest of the nodes.

Besides, we have guided our design to build the middleware over a REST interface.
Using that interface, each node exposes the semantic resources, i.e. RDF triples grouped in RDF graphs, it manages.
This enables the consumption of these contents by other semantic \ac{wot} solutions or REST services, even if they are not part of our middleware.

\section{Contributions}

% De una charla con Buján recuerdo que dijo que para O. Corcho...
%   + Contribuciones científicas: lo que nadie había hecho antes (innovación)
%   + Contribuciones técnicas: implementar o aplicar en otro campo un aspecto innovado por otro

% Contributions (por si necesito ideas)
% \item Analysis of strengths and weaknesses of both \ac{wot} and \ac{tsc} and comparison of both approaches.
% \item Design and implementation of a RESTful interface which makes possible to build a \ac{tsc} on top of it. % TODO revisar y poner todos los peros que haga falta para no traicionar la idea de Fielding
%       This resource-oriented interface enables the access to the the semantic knowledge provided by each device.
%       By doing so, we keep our middleware compliant with the \ac{wot} idea. % o integrable con otros servicios, o...
%       % es como explicar una dualidad: puedes usar nuestro middleware usando las primitivas de TSC, pero también puedes usarlo simplemente porque te viene bien para WoT semántico
% \item A search-aware architecture which promotes the end-to-end communication between devices. % and balances their load by administrative tasks according to their capacities.
%       This architecture tries to minimize the requests devices have to handle by enhancing their search mechanism.
%       The extra-tasks introduced by this enhancement are performed by nodes chosen according to their capacities.
% \item Design and comparison of two approaches which enable to act on the environment by writing in the fully distributed space.
% \item Exploration of the use of semantic content in several current embedded platforms.
%       We have tested several platforms, semantic frameworks and reasoners to verify the feasibility of the ideas presented in the dissertation.
% \item Design and implementation of most of the ideas presented in a middleware publicly available called \emph{Otsopack}. % al primo no le gustará, pero insisto
%       This middleware has been used in several projects for a range of scenarios.


\section{Relevant Publications}

% International JCR Journals
% International conferences
% Workshops % o si queda muy corto, incluirlo directamente en conclusions
\section{Future Work}

% Limitations => Future Work
\section{Final Remarks}



% ----------------------------------------------------------------------

%: ----------------------- On how to actuate through TSC ------------------------

% this file is called up by thesis.tex
% content in this file will be fed into the main document

%: ----------------------- introduction file header -----------------------
\begin{savequote}[50mm]
I tell you what, I speak slow so those of you with a Phd can understand.
\qauthor{Doug \textit{(Déjà Vu)}}
\end{savequote}


\chapter{Conclusions}
\label{cha:conclusions}
\newcommand{\pathchapseven}{7_conclusion}


% the code below specifies where the figures are stored
\ifpdf
    \graphicspath{{\pathchapseven/figures/PNG/}{\pathchapseven/figures/PDF/}{\pathchapseven/figures/}}
\else
    \graphicspath{{\pathchapseven/figures/EPS/}{\pathchapseven/figures/}}
\fi


%------------------------------------------------------------------------- 

% Presentar de qué ha ido la tesis
% 3 ptos con los que quiero que el lector se quede: searching, acting y compatibility with WoT

This dissertation explores the design of a completely distributed \ac{tsc} middleware.
In this middleware, each node is responsible of its own information, but together with other nodes they create a virtual space. % virtual, concept
Writing and reading from this space, the nodes communicate in a much more decoupled indirect style.

The distribution of the information avoids the dependency of the participants on other devices, but arises other challenges.
We have tackled two main challenges: how to search in a decentralized way and how to act on the environment through the space.
Besides, after an analysis we acknowledged the importance and benefits that the upcoming \ac{wot} has brought to the \acl{ubicomp}.
This has lead us to have make our final solution compatible with the \ac{wot} and other REST services.

To enhance the search process between devices limiting their dependencies in others we designed an architecture.
This architecture dynamically chooses an intermediary according to its capabilities to support the most limited devices.
The final result is that less requests are required in overall, specifically reducing the load on resource constrained devices.

% TODO refinar cuando esté hecha esa parte
We also presented two different acting mechanisms: X and Y.
Comparing them we reach to the conclusion that X proposes a less decoupled design, which in turn depends on a subscription mechanism.
This subscription mechanism is particularly challenging on dynamic environments.
The Y solution is more straightforward and more based in REST services consumption.
In contrast, it depends on reasoners and on some additional information about the services offered by the rest of the nodes.

Besides, we have guided our design to build the middleware over a REST interface.
Using that interface, each node exposes the semantic resources, i.e. RDF triples grouped in RDF graphs, it manages.
This enables the consumption of these contents by other semantic \ac{wot} solutions or REST services, even if they are not part of our middleware.

\section{Contributions}

% De una charla con Buján recuerdo que dijo que para O. Corcho...
%   + Contribuciones científicas: lo que nadie había hecho antes (innovación)
%   + Contribuciones técnicas: implementar o aplicar en otro campo un aspecto innovado por otro

% Contributions (por si necesito ideas)
% \item Analysis of strengths and weaknesses of both \ac{wot} and \ac{tsc} and comparison of both approaches.
% \item Design and implementation of a RESTful interface which makes possible to build a \ac{tsc} on top of it. % TODO revisar y poner todos los peros que haga falta para no traicionar la idea de Fielding
%       This resource-oriented interface enables the access to the the semantic knowledge provided by each device.
%       By doing so, we keep our middleware compliant with the \ac{wot} idea. % o integrable con otros servicios, o...
%       % es como explicar una dualidad: puedes usar nuestro middleware usando las primitivas de TSC, pero también puedes usarlo simplemente porque te viene bien para WoT semántico
% \item A search-aware architecture which promotes the end-to-end communication between devices. % and balances their load by administrative tasks according to their capacities.
%       This architecture tries to minimize the requests devices have to handle by enhancing their search mechanism.
%       The extra-tasks introduced by this enhancement are performed by nodes chosen according to their capacities.
% \item Design and comparison of two approaches which enable to act on the environment by writing in the fully distributed space.
% \item Exploration of the use of semantic content in several current embedded platforms.
%       We have tested several platforms, semantic frameworks and reasoners to verify the feasibility of the ideas presented in the dissertation.
% \item Design and implementation of most of the ideas presented in a middleware publicly available called \emph{Otsopack}. % al primo no le gustará, pero insisto
%       This middleware has been used in several projects for a range of scenarios.


\section{Relevant Publications}

% International JCR Journals
% International conferences
% Workshops % o si queda muy corto, incluirlo directamente en conclusions
\section{Future Work}

% Limitations => Future Work
\section{Final Remarks}



% ----------------------------------------------------------------------

%: ----------------------- conclusion ------------------------

% this file is called up by thesis.tex
% content in this file will be fed into the main document

%: ----------------------- introduction file header -----------------------
\begin{savequote}[50mm]
I tell you what, I speak slow so those of you with a Phd can understand.
\qauthor{Doug \textit{(Déjà Vu)}}
\end{savequote}


\chapter{Conclusions}
\label{cha:conclusions}
\newcommand{\pathchapseven}{7_conclusion}


% the code below specifies where the figures are stored
\ifpdf
    \graphicspath{{\pathchapseven/figures/PNG/}{\pathchapseven/figures/PDF/}{\pathchapseven/figures/}}
\else
    \graphicspath{{\pathchapseven/figures/EPS/}{\pathchapseven/figures/}}
\fi


%------------------------------------------------------------------------- 

% Presentar de qué ha ido la tesis
% 3 ptos con los que quiero que el lector se quede: searching, acting y compatibility with WoT

This dissertation explores the design of a completely distributed \ac{tsc} middleware.
In this middleware, each node is responsible of its own information, but together with other nodes they create a virtual space. % virtual, concept
Writing and reading from this space, the nodes communicate in a much more decoupled indirect style.

The distribution of the information avoids the dependency of the participants on other devices, but arises other challenges.
We have tackled two main challenges: how to search in a decentralized way and how to act on the environment through the space.
Besides, after an analysis we acknowledged the importance and benefits that the upcoming \ac{wot} has brought to the \acl{ubicomp}.
This has lead us to have make our final solution compatible with the \ac{wot} and other REST services.

To enhance the search process between devices limiting their dependencies in others we designed an architecture.
This architecture dynamically chooses an intermediary according to its capabilities to support the most limited devices.
The final result is that less requests are required in overall, specifically reducing the load on resource constrained devices.

% TODO refinar cuando esté hecha esa parte
We also presented two different acting mechanisms: X and Y.
Comparing them we reach to the conclusion that X proposes a less decoupled design, which in turn depends on a subscription mechanism.
This subscription mechanism is particularly challenging on dynamic environments.
The Y solution is more straightforward and more based in REST services consumption.
In contrast, it depends on reasoners and on some additional information about the services offered by the rest of the nodes.

Besides, we have guided our design to build the middleware over a REST interface.
Using that interface, each node exposes the semantic resources, i.e. RDF triples grouped in RDF graphs, it manages.
This enables the consumption of these contents by other semantic \ac{wot} solutions or REST services, even if they are not part of our middleware.

\section{Contributions}

% De una charla con Buján recuerdo que dijo que para O. Corcho...
%   + Contribuciones científicas: lo que nadie había hecho antes (innovación)
%   + Contribuciones técnicas: implementar o aplicar en otro campo un aspecto innovado por otro

% Contributions (por si necesito ideas)
% \item Analysis of strengths and weaknesses of both \ac{wot} and \ac{tsc} and comparison of both approaches.
% \item Design and implementation of a RESTful interface which makes possible to build a \ac{tsc} on top of it. % TODO revisar y poner todos los peros que haga falta para no traicionar la idea de Fielding
%       This resource-oriented interface enables the access to the the semantic knowledge provided by each device.
%       By doing so, we keep our middleware compliant with the \ac{wot} idea. % o integrable con otros servicios, o...
%       % es como explicar una dualidad: puedes usar nuestro middleware usando las primitivas de TSC, pero también puedes usarlo simplemente porque te viene bien para WoT semántico
% \item A search-aware architecture which promotes the end-to-end communication between devices. % and balances their load by administrative tasks according to their capacities.
%       This architecture tries to minimize the requests devices have to handle by enhancing their search mechanism.
%       The extra-tasks introduced by this enhancement are performed by nodes chosen according to their capacities.
% \item Design and comparison of two approaches which enable to act on the environment by writing in the fully distributed space.
% \item Exploration of the use of semantic content in several current embedded platforms.
%       We have tested several platforms, semantic frameworks and reasoners to verify the feasibility of the ideas presented in the dissertation.
% \item Design and implementation of most of the ideas presented in a middleware publicly available called \emph{Otsopack}. % al primo no le gustará, pero insisto
%       This middleware has been used in several projects for a range of scenarios.


\section{Relevant Publications}

% International JCR Journals
% International conferences
% Workshops % o si queda muy corto, incluirlo directamente en conclusions
\section{Future Work}

% Limitations => Future Work
\section{Final Remarks}



% ----------------------------------------------------------------------


% --------------------------------------------------------------
%:                  BACK MATTER: appendices, refs,..
% --------------------------------------------------------------

% the back matter: appendix and references close the thesis
\backmatter


%: ----------------------- appendix ------------------------

%\appendix


%: ----------------------- bibliography ------------------------



% The section below defines how references are listed and formatted
% The default below is 2 columns, small font, complete author names.
% Entries are also linked back to the page number in the text and to external URL if provided in the BibTex file.


% Original version:

% PhDbiblio-url2 = names small caps, title bold & hyperlinked, link to page 
%\begin{multicols}{2} % \begin{multicols}{ # columns}[ header text][ space]
%\begin{tiny} % tiny(5) < scriptsize(7) < footnotesize(8) < small (9)
%
%\bibliographystyle{Latex/Classes/PhDbiblio-url2} % Title is link if provided
%\renewcommand{\bibname}{References} % changes the header; default: Bibliography
%
%\bibliography{9_backmatter/references} % adjust this to fit your BibTex file
%
%\end{tiny}
%\end{multicols}



% Show all bibliography entries
%\nocite*



% If we want bibliography backreference, use unsrt first and the desidered one after

%\bibliographystyle{unsrt} % Defines the bibliography style

%\bibliographystyle{alpha} % Defines the bibliography style

%\bibliographystyle{apa-good} % Defines the bibliography style
%\bibliographystyle{natbib} % Defines the bibliography style

%\bibliographystyle{plainurl}

%\renewcommand{\bibname}{References} % changes the header; default: Bibliography

%To include the references/works cited/bibliography in your Table of Contents, right before the bibliography command, use the command
%\addcontentsline{toc}{section}{References}
%

\bibliography{9_backmatter/ref_rest,9_backmatter/ref_ubicomp,9_backmatter/ref_space,9_backmatter/ref_semantic,9_backmatter/ref_self,9_backmatter/ref_others} % adjust this to fit your BibTex file

% --------------------------------------------------------------
% Various bibliography styles exit. Replace above style as desired.

% in-text refs: (1) (1; 2)
% ref list: alphabetical; author(s) in small caps; initials last name; page(s)
%\bibliographystyle{Latex/Classes/PhDbiblio-case} % title forced lower case
%\bibliographystyle{Latex/Classes/PhDbiblio-bold} % title as in bibtex but bold
%\bibliographystyle{Latex/Classes/PhDbiblio-url} % bold + www link if provided

%\bibliographystyle{Latex/Classes/jmb} % calls style file jmb.bst
% in-text refs: author (year) without brackets
% ref list: alphabetical; author(s) in normal font; last name, initials; page(s)

% Other styles such as alpha don't go well with natbib, and I want to use \citep, \citeauthor, etc.!
\bibliographystyle{plainnat} % calls style file plainnat.bst
% in-text refs: author (year) without brackets
% (this works with package natbib)


% --------------------------------------------------------------


% HOWTO Citar páginas concretas o comentarios (blogs):
% http://tex.stackexchange.com/questions/89143/citation-with-some-text-inside-square-brackets

%: Declaration of originality

% this file is called up by thesis.tex
% content in this file will be fed into the main document

%: ----------------------- introduction file header -----------------------


% the code below specifies where the figures are stored
\ifpdf
    \graphicspath{{9_backmatter/figures/PNG/}{9_backmatter/figures/PDF/}{9_backmatter/figures/}}
\else
    \graphicspath{{9_backmatter/figures/EPS/}{9_backmatter/figures/}}
\fi

% Thesis statement of originality -------------------------------------

% Depending on the regulations of your faculty you may need a declaration like the one below. This specific one is from the medical faculty of the university of Dresden.

\begin{declaration}        %this creates the heading for the declaration page

I herewith declare that I have produced this work without the prohibited assistance of third parties and without making use of aids other than those specified; notions taken over directly or indirectly from other sources have been identified as such. This work has not previously been presented in identical or similar form to any examination board.

The dissertation work was conducted from 2010 to \the\year \ under the supervision of
\Advisor 
~and~
\Advisortwo
~at 
the University of Deusto.

\vspace{10mm}

% City
Bilbao,

% Signature figure

%\begin{flushright}
%\begin{figure}[htbp!]
%\end{figure}
%\includegraphics{signature}%
%\end{flushright}


\end{declaration}



%-------------------------------------------------------------------------

\thispagestyle{empty}

\hfill
\vfill
\medskip

\begin{center}
\noindent
This dissertation was finished writing 
in 
Bilbao
on
% Date
\today 
\end{center}

\vfill
\medskip
\vspace{1cm}
\bigskip


%-------------------------------------------------------------------------

% Blank page
\clearpage
\thispagestyle{empty}

\hfill
\vfill
\medskip

\begin{center}
\textit{This page is intentionally left blank}
\end{center}

\vfill
\medskip
\vspace{1cm}
\bigskip


% ----------------------------------------------------------------------





\end{document}
